\section{Conditional pencil contributions}
\label{app:pencil}
For this calculation impose the additional condition that every member
of both pencils is reduced and that the displayed toric members are
their only reducible members. This is an explicit hypothesis on the
pencils, not a consequence of the cubic form; a complete genericity
check for the anticanonical family is not supplied here. The isolated sheaves of Section~\ref{subsec:x9-moving-probes}
do not depend on this hypothesis.

For $n=0,1$, consider the sheaves $\cO_T$ and $I_{p,T}$, respectively,
on members $T\in|P_i|$, pushed forward to $X$. Use the ample
polarization $J_*=(6,18,16,-11,-8,7)$ in the basis
\eqref{eq:x9-divisor-basis}; it is the sum of the six primitive ambient
nef generators. These sheaves form open and closed smooth components
of the stable pure-sheaf moduli, with contributions
\begin{equation}
 \begin{array}{c|cc|cc}
 &\mathcal M_i(0)&\mathcal M_i(1)&
 \Omega_i^{\rm seed}(0)&\Omega_i^{\rm seed}(1)\\\hline
 i=5&\mathbb P^1&X&-2&232\\
 i=6&\mathbb P^1&\operatorname{Bl}_{C_1}X&-2&230 .
 \end{array}
 \label{eq:x9-moving-seeds}
\end{equation}
The superscript denotes only these specified components, not the
full index at the same charge.

Here is a component and stability argument, including the reducible
members. For $T=A+B$, a proper subunion $A$ supports the maximal
subsheaf $\cO_A(-B)\subset\cO_T$. A subsheaf of $I_{p,T}$ supported
on $A$ has no larger reduced Hilbert polynomial. The strict comparison
of the linear coefficients with that of $\cO_T$ is
\begin{equation}
 \delta_T(A)=
 [J_*\cdot A(A+2B)](J_*^2\cdot T)
 -(J_*\cdot T^2)(J_*^2\cdot A)>0.
 \label{eq:x9-probe-stability}
\end{equation}
For the ordered components in \eqref{eq:x9-probe-pencils}, and subunions
$1,2,3,12,13,23$, the numerators are
\begin{equation}
 \begin{array}{c|rrrrrr}
 &1&2&3&12&13&23\\\hline
 P_5&2952&3690&1230&1230&3690&2952\\
 P_6&3740&4876&1482&1488&4628&3982 .
 \end{array}
 \label{eq:x9-probe-stability-data}
\end{equation}
Subsheaves of full multirank have a nonzero lower-dimensional quotient.
The remaining pencil members are integral by hypothesis, where rank-one
stability is automatic. Thus all the specified sheaves are stable,
including those with the point on a double curve.

The smooth member determines the line-bundle cohomology:
$h^0(X,\cO_X(P_i))=2$, with no higher cohomology.
The divisor sequence gives $H^1(T,\cO_T)=0$ for every member. Consequently
$\operatorname{Ext}^1_X(\cO_T,\cO_T)$ is the one-dimensional pencil
tangent space. For a point $p$ outside the base locus, the evaluation
map on the two sections is surjective, so
$R\operatorname{Hom}(\cO_X(-P_i),I_p)=\CC$ in degree zero.
The sequence
\begin{equation}
 0\longrightarrow\cO_X(-P_i)\longrightarrow I_p
 \longrightarrow I_{p,T}\longrightarrow0
 \label{eq:x9-probe-twist}
\end{equation}
identifies $I_{p,T}$ with the spherical twist of $I_p$ about
$\cO_X(-P_i)$. This twist is an autoequivalence
\cite[Theorem~1.2]{Seidel:2001Twists}; it preserves the
three-dimensional $\operatorname{Ext}^1(I_p,I_p)$.
For $P_6$ and $p\in C_1$, every pencil member is smooth along the
base curve. The normal-deformation sequence then bounds
$\dim\operatorname{Ext}^1_X(I_{p,T},I_{p,T})$ by
$2+h^0(K_T)=3$. The smooth three-dimensional incidence variety
$\operatorname{Bl}_{C_1}X$ already supplies these deformations.
Thus the natural proper parameter spaces are open and closed smooth
components also at the base locus and singular fibers.

Their Behrend functions are $-1$. Since
$\chi(X)=2(6-122)=-232$ and
$\chi(\operatorname{Bl}_{C_1}X)=\chi(X)+\chi(C_1)=-230$,
the entries in \eqref{eq:x9-moving-seeds} follow.
In the charge convention of \eqref{eq:x9-supported-mukai}, the two
families have
\begin{equation}
 \Gamma_{5,n}=(0,P_5,0,1-n),\qquad
 \Gamma_{6,n}=(0,P_6,-\tfrac12 C_1,\tfrac{11}{12}-n).
 \label{eq:x9-points-seed-charges}
\end{equation}
These results fix neither the other components nor the physical
stability condition near the cooperative contraction.
